\paragraph{\S102. Abelianizators, integrators and the dual abelian pairing (pp.\ 378--381).}
For a small category $A$ and $B=A^o$, the section returns explicitly to \S93 and rebuilds the relation between integrators and abelianization through the dual pair of presheaf categories. Comparing (2) with \S99, (3), p.\ 362, gives the finite additive-envelope equivalence $\mathrm{Add}(A^o)\simeq\mathrm{Add}(A)^o$. The proposed infinite analogue (3') is entirely canceled, together with the words introducing it; it is therefore not an assertion that $\mathrm{Addinf}$ commutes with passage to the opposite category. The infinite envelope still enters separately in (4), because quasi-special integrators exist even when special ones need not. The revised pairing (*) between $A^\wedge_{ab}$ and $\mathrm{Addinf}(B)$ is biadditive and exact in its first variable; (5) is correspondingly revised from additive to exact functors. Extension to (6), with second argument in all of $B^\wedge_{ab}$, preserves small inductive limits in both variables (p.\ 378). A complex $L_\bullet$ in $\mathrm{Addinf}(B)$ consequently gives the exact complex-valued functor (7). An abelianizator for $A$ is then a pair $(L_\bullet,\lambda)$ supplying the comparison in (8) between the abelianized presheaf construction and absolute abelianization of homotopy types. The latter is interpreted through the integrators of \S100 on an arbitrary modelizing category, avoiding a privileged auxiliary simplicial category; the earlier constructions of \S\S92--93 are thus placed in a more general functorial setting (p.\ 379). The aim is simultaneous chain-level computation for varying $X$, with a formula extending from $\mathrm{Wh}_A(X)$ to arbitrary abelian presheaves. A quasi-special integrator on $B$, namely a resolution $L_\bullet\to\underline Z_B$ in (10), supplies an abelianizator for $A$ and its canonical comparison. The source explicitly says \emph{for A} here, where the witness substitutes B. The converse is proposed as an equivalence, perhaps even an isomorphism, but is not pursued. Enlarging from special and quasi-special resolutions to arbitrary projective resolutions motivates extending the definition of abelianizator through the whole pairing (6), while existence and uniqueness up to chain homotopism are inherited from the resolution construction (p.\ 380). Formula (11) identifies $B^\wedge_{ab}$ with the cocontinuous abelian-group-valued functors on $A^\wedge_{ab}$. Projective objects give exact functors; the converse is explicitly left to be checked. The image of $\underline Z_B$ is the inductive-limit functor in (13), equivalently the integration functor $p^{ab}_{A!}$ for $p_A:A\to\Delta_0$ in (14), so (15) describes standard abelianizators by projective resolutions of that functor. This makes precise an approach that the cellular and test-category questions of \S91 had left uncertain (p.\ 381). The tentative converse on p.\ 381 must be read together with its explicit rejection in \S106, p.\ 410: already for the final category, exactness of tensor product detects flatness, or torsion-freeness over $\underline Z$, rather than projectivity or freeness. Section~103 revisits the proof strategy and changes the terminology from abelianizator toward integrator and cointegrator.

\paragraph{\S103. Cointegrators, integrators and a change of approach (pp.\ 382--383).}
Returning on 11 August after three weeks away from the notes, the author reconsiders the proposed proof that an integrator on the opposite category gives an abelianizator. His earlier route used Quillen's comparison of a category with its opposite: the functor $(A_{/X})^o\to A^o=B$ is a cofibration with discrete fibers, so its exact direct image gives the displayed comparison of derived sections. The source retains an upper star on the first coefficient category in that display, unlike the arXiv witness, and has a handwritten margin VII beside the still-typed reference to part VI. This cohomological route is then rejected as an awkward way to establish a result about homology. The proposed natural proof is postponed until an overall review makes the relations among homology, cohomology and abelianization conceptually clear (pp.\ 382--383). The terminology is revised at the same time, explicitly in light of \S69: a projective resolution of the constant abelian presheaf $\underline Z_A$ computes $R\Gamma(A,-)$, hence cointegration, and should be called a \emph{cointegrator for A}. For the opposite category $B=A^o$, the same complex computes homology and is an \emph{integrator for B}. Such complexes also give the abelianizators under discussion in \S\S93 and 102. A full converse remains only a possibility, but the existence and uniqueness up to homotopism of the resolution construction make other putative abelianizators seem dispensable. The author therefore proposes dropping the term abelianizator in favor of integrator, paired with cointegrator. The paragraph literally calls the object an abelianizator for B while still writing $X\in A^\wedge$ and $\mathrm{Wh}_A(X)$ in the associated computation; this surviving mismatch of base indices is preserved in the text. The new terminology and the decision to rebuild the formalism lead directly to \S104 (p.\ 383).

\paragraph{\S104. Pseudotopoi, abelianization\newline and internal and external operations (pp.\ 384--393).}
The conceptual review promised in \S103 begins with the full subcategories $\underline{\mathrm{Hom}}_!$ and $\underline{\mathrm{Hom}}^!$ of ordinary functor categories, preserving small inductive and projective limits respectively. Proposition~1 assumes a category $\mathcal A$ with small inductive limits and a small full subcategory generating with respect to monomorphisms; citing SGA~4 I, 8.12.7, it identifies $\mathcal A$ with the projective-limit-preserving functors from $\mathcal A^o$ to sets and deduces small projective limits. The defining test literally refers to objects $Z$ of $C$ after naming the generating subcategory $A$; this index mismatch remains in the source (p.\ 384). Such a category is called a pseudo-topos. Corollary~2 gives the corresponding adjoint criterion and the duality between the two Hom categories. The handwritten revision of (4) places an outer opposite on the first Hom category and cancels the one on the second; (6) explicitly expresses the ordinary reversal of natural transformations under passage to opposite categories. This convention matters for the later formula (37), whose missing opposite is discussed in \S105. Both sides of (4) are identified with the separately projective-limit-preserving bifunctors (5). For actual topoi these are valued sheaves, and for presenting sites they become bisheaves; their category is described as the $2$-product of the two topoi (pp.\ 385--386). Part B introduces $\mathcal A_{ab}$ and its forgetful functor, with left adjoint $\mathrm{Wh}_{\mathcal A}$ when it exists. On p.\ 386 the source literally says that this forgetful functor commutes with ``small direct limits,'' although the displayed adjunction argument and the preceding corollary require the projective-limit statement; no mathematical correction is made in the text. For a topos, the abelian category of abelian group objects has exact small filtered inductive limits, a small generating subcategory, enough injectives and a small cogenerating subcategory, so it is both a pseudotopos and a pseudocotopos. More generally, the pseudocotopos condition on $\mathcal A_{ab}$ ensures existence of abelianization. Proposition~2 and its corollary give the universal property (11)--(12): precomposition with $\mathrm{Wh}_{\mathcal A}$ identifies cocontinuous functors from $\mathcal A_{ab}$ to an additive pseudocotopos with those from $\mathcal A$. The dual formula (13), followed by the choice of abelian groups as coefficients, yields (14), interpreting abelian group objects as abelian-group-valued sheaves. The witness drops the essential co in two pseudocotopos hypotheses on p.\ 388 (pp.\ 386--388). Part C defines tensor product by representability of bilinear maps and recalls its construction from abelianization and cokernels; those ingredients can suffice without the full pseudotopos assumption. The canonical comparison (17), associativity, commutativity and the unit $\underline Z_{\mathcal A}=\mathrm{Wh}_{\mathcal A}(e)$ follow. The explanation of the final object on p.\ 389 again literally invokes ``small direct limits,'' where the preceding theory supplies projective limits. The notation (19) distinguishes $\mathcal A^{\mathcal M}=\underline{\mathrm{Hom}}_!(\mathcal A,\mathcal M)$, called cosheaves, from $\mathcal A_{\mathcal N}=\underline{\mathrm{Hom}}^!(\mathcal A^o,\mathcal N)$, called sheaves. The handwritten condition that $\mathrm{Wh}_{\mathcal A}$ exist is retained; the witness omits it at this point. Under $\mathcal N=\mathcal M^o$, (20) identifies the two categories only after taking an opposite. A later change from $A^\wedge$ to the opposite presheaf topos $B^\wedge$, $B=A^o$, is promised to remove this variance reversal while keeping the coefficient category fixed. The sentence describing the base of the sheaves on p.\ 390 actually names $\mathcal N$ as the topos instead of $\mathcal A$; the scan settles that literal letter (pp.\ 389--390). To define the external action on cosheaves, tensoring with a fixed abelian sheaf must preserve small inductive limits. Formula (22) then sets $(F\circledast L)(G)=L(F\otimes G)$, and (23) abbreviates evaluation as $H*L=L(H)$; the circled star and plain star have distinct roles. The required preservation is reformulated through representability of the presheaf-valued internal Hom (24) and the Cartan relation (25), leading to the stated assumptions that tensor products, corresponding internal Hom objects and abelianization exist. The source writes F, H after (24), where the witness has F, G; it also cites (15) for abelianization, whose definition is (10), and briefly says ``sheaves'' and ``presheaves'' where the surrounding discussion distinguishes sheaves and cosheaves. Under these hypotheses the action (26) has associativity (27) and unit (28) (pp.\ 390--392). Duality gives the exterior Hom action (29)--(30), characterized by $\underline{\mathrm{Hom}}(F,K)(G)=K(G\otimes F)$, while the un-underlined Hom of (31) again denotes evaluation. The source uses K in (30) and (30'), whereas the witness substitutes L in the inner Hom. With $\mathcal N=(Ab)$, these operations recover ordinary Hom and internal Hom of abelian sheaves. An analogous interpretation of the cosheaf tensor operation is not established for a general pseudotopos; the last line of p.\ 392 asks for a choice of $\mathcal N$ despite discussing $\mathcal M$-valued cosheaves, and p.\ 393 returns to $\mathcal M$ without explicitly reassigning the symbol. In the presheaf case $\mathcal A=A^\wedge$, the promised duality identifies $\mathcal A^{\mathcal M}$ with sheaves on $B^\wedge$, and with $B^\wedge_{ab}$ for $\mathcal M=(Ab)$. For sets, cocontinuous endofunctors of abelian groups are tensoring with an abelian group, so both exterior operations recover the familiar tensor product. Compatibility when $A=A^o$ literally, as for sums of commutative one-object categories, is left unchecked. The closing distinction between object-valued and sheaf-valued Hom, and between object-valued star and cosheaf-valued circled star, prepares the systematic presheaf duality of \S105 and the computational formulaire of \S106 (pp.\ 392--393).
